\section{Primitive-recursive Catalan bijections}
\label{sec:combinatorics}

This section proves the principal counting interpretations listed on the two
source pages.  Each map is an algorithm on bounded words or lists, and each
inverse is described.  Thus the arguments are not merely classical counting
proofs; they directly instantiate Lemma~\ref{lem:finite-bijection}.

\subsection{Words, parentheses, lattice paths, and mountains}

The Dyck-word interpretation is Theorem~\ref{thm:reflection}.  Replacing $U$
by an opening parenthesis and $D$ by a closing parenthesis gives correctly
matched parenthesis strings: the height is the number of currently open
parentheses.  Replacing $U,D$ by right and up steps gives monotone paths from
$(0,0)$ to $(n,n)$ that remain on the chosen side of the diagonal.  Replacing
them by slopes $+1,-1$ gives mountains that do not fall below the starting
level.  These are literal letter substitutions, so their inverses are
immediate.

The column-height representation on Wikipedia is also direct.  For the
$j$th right step, record the number $h_j$ of up steps preceding it.  A path
below the diagonal gives
\[
 0\le h_1\le h_2\le\cdots\le h_n,
 \qquad h_j\le j-1.
\tag{3.1}\label{eq:height-sequence}
\]
Conversely, insert the required number of up steps before each right step and
then the remaining up steps.  This gives a primitive-recursive bijection with
nondecreasing sequences satisfying Equation~\eqref{eq:height-sequence}.

\subsection{Full binary trees and complete parenthesizations}

A full ordered binary tree is either a leaf or an internal node with an ordered
left and right subtree.  Define its Dyck code recursively by
\[
 \operatorname{code}(\text{leaf})=\varepsilon,
 \qquad
 \operatorname{code}(\operatorname{node}(L,R))
   =U\operatorname{code}(L)D\operatorname{code}(R).              \tag{3.2}\label{eq:tree-code}
\]
If the tree has $n$ internal nodes, the word has $n$ $U$'s and $n$ $D$'s.
The first-return decomposition in Theorem~\ref{thm:segner} recovers $L$ and
$R$, so Equation~\eqref{eq:tree-code} is a bijection.  Hence $\Cat_n$ counts full
ordered binary trees with $n$ internal vertices and $n+1$ leaves.

A complete parenthesization of $n+1$ ordered factors has a unique outermost
binary operation; its two operands are parenthesized initial and final blocks.
Make this operation the root and recurse.  Reading a tree back as an operation
gives the inverse.  The same proof covers binary parse trees, bracketings of
ordered exponentials, and any fixed binary operator.  Adding a plant edge and
suppressing leaves converts a full binary tree into the planted planar
trivalent tree described by MathWorld.

\subsection{Rooted plane trees}

For a rooted ordered tree with $n$ edges, perform its contour traversal.  Write
$U$ when an edge is traversed away from the root for the first time and $D$
when it is traversed back.  The current depth is the Dyck height, so the output
is a Dyck word of semilength $n$.  Conversely, scan a Dyck word with a stack:
on $U$, create the next child of the current vertex and push it; on $D$, pop.
The prefix condition prevents underflow and the final height empties the
stack.  These two routines are inverse, proving that $\Cat_n$ counts rooted
plane trees with $n$ edges, equivalently $n+1$ vertices.

The familiar left-child/right-sibling representation gives the same
bijection through binary trees: the left pointer is the first child and the
right pointer is the next sibling.  This alternative is useful in Lean because
both trees can share one bounded-array representation.

\subsection{Triangulations of a convex polygon}

Fix a distinguished boundary edge of a convex $(n+2)$-gon.  In any
triangulation, exactly one triangle is incident with that edge.  If its third
vertex leaves $i+2$ vertices on one side and $n-i+2$ on the other, the remaining
diagonals are independent triangulations of those two polygons.  Thus the
numbers $T_n$ satisfy
\[
 T_0=1,\qquad T_{n+1}=\sum_{i=0}^{n}T_iT_{n-i}.                  \tag{3.3}
\]
More strongly, the incident triangle is a computable root, and recursively
placing the two subtriangulations to its left and right is a bijection with
full ordered binary trees.  Therefore $T_n=\Cat_n$.  The same recursion proves
Euler's polygon-division problem and Segner's formula
$E_m=\Cat_{m-2}$ for $m\ge2$.

\subsection{Noncrossing handshakes and pair partitions}

Place $2n$ labeled points around a circle.  In a noncrossing perfect matching,
point $1$ is paired with an even point $2i+2$.  The $2i$ points inside that
chord and the remaining $2(n-1-i)$ points outside it are matched
independently.  This is an explicit first-chord decomposition with inverse,
so the count satisfies Segner's recurrence and is $\Cat_n$.

Equivalently, scan the points in cyclic order and write $U$ at the first
endpoint and $D$ at the second endpoint of each chord, where chords are
processed by nesting.  Noncrossing is exactly the stack condition.  This also
proves the Wikipedia statement about noncrossing partitions of $[2n]$ whose
blocks all have size two.

\subsection{Noncrossing set partitions and the Bell bound}

Encode a partition of $[n]$ by its restricted-growth word.  It is noncrossing
when there do not exist $a<b<c<d$ with $a,c$ in one block and $b,d$ in another.
This is a bounded decidable predicate.

Let $P(x)$ be the formal counting series for noncrossing partitions.  In a
nonempty partition, list the elements of the block containing $1$ as
$1=b_0<b_1<\cdots<b_r$.  The $r+1$ cyclic intervals after these elements carry
independent noncrossing partitions.  Marking the $r+1$ elements of the root
block gives, coefficientwise,
\[
 P=1+\sum_{r\ge0}x^{r+1}P^{r+1}
  =1+\frac{xP}{1-xP}.                                           \tag{3.4}\label{eq:NCseries}
\]
Multiplication by $1-xP$ reduces this to $P=1+xP^2$.  By formal uniqueness,
$P=C$, so there are $\Cat_n$ noncrossing partitions of $[n]$.  The decomposition
and reconstruction are bounded list algorithms; the geometric terminology is
not essential to the proof.

The inclusion of noncrossing partitions into all set partitions is injective,
so
\[
 \Cat_n\le \operatorname{Bell}(n).                               \tag{3.5}\label{eq:bell-bound}
\]
This is an \ipra{} inequality between bounded counts, not an appeal to cardinal
comparison of abstract sets.

\subsection{Two-row standard Young tableaux}

Given a standard Young tableau of rectangular shape $(n,n)$, scan the numbers
$1,\ldots,2n$ and write $U$ if the number is in the first row and $D$ if it is
in the second.  Each row contains $n$ entries.  The $j$th first-row entry is
smaller than the $j$th second-row entry exactly when every prefix has at least
as many $U$'s as $D$'s.  Thus the word is Dyck.

Conversely, place the positions of the $U$'s increasingly in the first row and
the positions of the $D$'s increasingly in the second.  The Dyck prefix
condition is precisely the column inequality.  This proves that the number of
$2$-by-$n$ standard tableaux is $\Cat_n$.  Evaluating the hook-length product
also gives
\[
 \frac{(2n)!}{n!(n+1)!}=\Cat_n,
\]
so the hook-length specialization agrees with the bijection.

\subsection{The restricted sequences on Wikipedia}

Let $a_i$ be the height immediately after the $i$th $U$ in a Dyck word.  Then
\[
 a_1=1,
 \qquad 1\le a_{i+1}\le a_i+1.                                  \tag{3.6}\label{eq:restricted-seq}
\]
Thus the sequence may rise by at most one and may fall by any amount while
remaining positive.  Conversely, from such a sequence insert
$a_i-a_{i+1}+1$ letters $D$ between the $i$th and $(i+1)$st $U$, and insert
$a_n$ letters $D$ at the end.  These numbers are nonnegative and telescope to
$n$ down steps.  The construction stays nonnegative and inverts height
recording.  Hence these length-$n$ sequences are counted by $\Cat_n$.

\subsection{Pattern-avoiding and stack-sortable permutations}

A permutation avoids $231$ precisely when, after writing it as $u\,n\,v$ at
its largest entry, every entry of $u$ is smaller than every entry of $v$, and
both standardized subwords avoid $231$.  If some $x\in u$ exceeds some
$y\in v$, then $x,n,y$ is a $231$ pattern; the converse is immediate by
locating the largest member of any alleged pattern.  Therefore the left word
uses the smallest $i$ values for some $i$, the right word uses the remaining
values, and
\[
 A_{n+1}=\sum_{i=0}^{n}A_iA_{n-i},\qquad A_0=1.                  \tag{3.7}
\]
Standardization and unstandardization are primitive recursive, so
$A_n=\Cat_n$ by an actual decomposition bijection.

Reverse, complement, and inverse are explicit involutions on permutations.
They show that $132,213,231,312$ have the same avoidance count.  For the
remaining orbit, the Robinson--Schensted row-insertion algorithm is primitive
recursive and its reverse row-deletion algorithm is an explicit inverse.  The
special Schensted lemma needed here also has a bounded inductive proof: the
row-bumping lemma says that insertion paths for successively smaller letters
move strictly left and strictly downward.  Hence a decreasing subsequence of
length $r$ forces $r$ distinct rows.  Conversely, tracing the bumping
predecessors of the leftmost entry in each successive row, and then reversing
the recorded insertion times, constructs a decreasing subsequence with one
entry from every occupied row.  Thus a permutation avoids $321$ exactly when
its insertion tableau has at most two rows.  Therefore its number is
\[
 \sum_{k=0}^{\floor{n/2}} f_{(n-k,k)}^2,
\]
where the two-row hook calculation gives
$f_{(n-k,k)}=\binom nk-\binom n{k-1}$.  Equation~\eqref{eq:square-sum} identifies
this sum with $\Cat_n$.  Reversal covers $123$.  Consequently every pattern of
length three has exactly $\Cat_n$ avoiding permutations.

The recursive stack-sorting operator satisfies
$S(u\,n\,v)=S(u)S(v)n$.  Its output is increasing if and only if $S(u)$ and
$S(v)$ are increasing and every value in $u$ is smaller than every value in
$v$.  By the preceding characterization this is equivalent, inductively, to
$231$-avoidance.  Thus the stack-sortable permutations are also counted by
$\Cat_n$.

\subsection{Staircase tilings and related tree models}

A tiling of the staircase Ferrers shape $(n,n-1,\ldots,1)$ by $n$ rectangles
has a unique rectangle containing the northwest corner.  Extending its
southeast corner to the antidiagonal splits the remaining cells into two
smaller staircase regions; empty regions are allowed.  Recursing produces a
full ordered binary tree with one internal node per rectangle.  Conversely,
place the root rectangle across the corresponding antidiagonal interval and
recurse on the two subregions.  The split coordinates are integral and
bounded, so this is a primitive-recursive inverse pair.  It proves the
staircase-tiling interpretation stated on Wikipedia.

MathWorld also lists admissible diagonals in positive integral frieze patterns
and abstract flexagon configurations.  Their published correspondences pass
through triangulations or ordered pairs of ``pats.''  The triangulation side
has been proved above.  A complete formal proof for those specialized object
formats would require fixing their exact source conventions; this package
records them as formalization targets rather than silently changing the
definitions.  See the coverage ledger in Section~\ref{sec:coverage}.
